% -*- coding: utf-8 -*-
% !TEX root = ../main.tex

\subsection{Pipeline expérimental commun}
\label{ssec:methode_pipeline}

Toutes les méthodes partagent les mêmes partitions de données, le même
tokenizer BPE, le même modèle Word2Vec W1 et la même initialisation de la tête
dense M1. Cette base commune est essentielle : si le backbone ou la tête
initiale changeait entre deux expériences, la différence de performance ne
pourrait plus être attribuée uniquement au critère de pruning.

Le modèle W1 est appris une fois avec l'objectif CBOW, puis figé. Pour chaque
texte, les embeddings d'entrée des unités BPE sont moyennés afin d'obtenir
$\boldsymbol{x}\in\mathbb{R}^{100}$. La tête est entraînée avec la
binary cross-entropy appliquée au logit $z_r$ :
\begin{equation}
    \mathcal{L}(\boldsymbol{\theta})
    =-\frac{1}{B}\sum_{r=1}^{B}
    \left[
        y_r\log\sigma(z_r)
        +(1-y_r)\log\bigl(1-\sigma(z_r)\bigr)
    \right].
    \label{eq:loss_bce_methodologie}
\end{equation}
Les variantes prunées utilisent exactement cette même fonction de coût.

Le protocole commun suit la séquence suivante :
\begin{enumerate}
    \item entraîner la tête dense M1 et conserver le meilleur état selon la
    loss de validation ;
    \item calculer les scores d'importance sur M1 ou sur le dernier modèle
    réentraîné ;
    \item construire un masque et mettre à zéro les poids sélectionnés pour
    obtenir M2 ;
    \item réentraîner les poids actifs sous contrainte du même masque afin
    d'obtenir M3 ;
    \item comparer M2 et M3 à M1 ainsi qu'à la petite architecture dense M4,
    entraînée directement avec un nombre de poids voisin de la cible.
\end{enumerate}

La configuration commune actuellement inscrite dans les scripts est résumée
dans le tableau~\ref{tab:hyperparametres_communs}.

\begin{table}[!ht]
    \centering
    \begin{tabularx}{\textwidth}{|l|X|}
        \hline
        Élément & Configuration actuelle \\ \hline
        Données & Jeu \texttt{Sp1786/multiclass-sentiment-analysis-dataset}, classe neutre retirée : $19\,583$ textes d'entraînement, $3\,277$ de validation et $3\,276$ de test. \\ \hline
        Tokenizer & BPE de vocabulaire maximal $|\mathcal{V}|=1\,000$, fréquence minimale 2. \\ \hline
        Word2Vec W1 & CBOW, dimension 100, fenêtre de deux unités de chaque côté, 5 négatifs par cible. \\ \hline
        Entraînement W1 & 5 époques, batch 512, Adam, taux d'apprentissage $10^{-3}$. \\ \hline
        Tête H1/M1 & $100\rightarrow1\,024\rightarrow1$. \\ \hline
        Tête H2/M1 & $100\rightarrow1\,024\rightarrow128\rightarrow1$. \\ \hline
        Entraînement des têtes & Au plus 30 époques, batch 256, Adam, taux $10^{-3}$, patience 4 sur la loss de validation. \\ \hline
        Fonction de coût & Binary cross-entropy calculée directement à partir du logit. \\ \hline
        Graine actuelle & 42 pour l'initialisation, les tirages et l'entraînement. \\ \hline
    \end{tabularx}
    \caption{Hyperparamètres communs au pipeline actuel.}
    \label{tab:hyperparametres_communs}
\end{table}

\subsection{Dérivation du score Optimal Brain Damage}
\label{ssec:methode_obd}

\subsubsection{Approximation locale au second ordre}

OBD cherche à estimer l'augmentation de la fonction de coût provoquée par la
suppression d'un paramètre \cite{lecun1989obd}. Soit
$\boldsymbol{\theta}$ le vecteur des paramètres de la tête entraînée et
$\boldsymbol{\delta}$ une perturbation de ces paramètres. Au voisinage de
$\boldsymbol{\theta}$, un développement de Taylor au second ordre donne
\begin{equation}
    \Delta\mathcal{L}
    =\mathcal{L}(\boldsymbol{\theta}+\boldsymbol{\delta})
     -\mathcal{L}(\boldsymbol{\theta})
    \simeq
    \boldsymbol{g}^{\mathsf{T}}\boldsymbol{\delta}
    +\frac{1}{2}
     \boldsymbol{\delta}^{\mathsf{T}}
     H\boldsymbol{\delta},
    \label{eq:taylor_obd}
\end{equation}
où $\boldsymbol{g}=\nabla_{\boldsymbol{\theta}}\mathcal{L}$ et
$H=\nabla_{\boldsymbol{\theta}}^2\mathcal{L}$ désignent respectivement le
gradient et la Hessienne.

Le passage de l'\autoref{eq:taylor_obd} au score OBD repose sur trois
hypothèses. Premièrement, le réseau est supposé proche d'un minimum, de sorte
que $\boldsymbol{g}\simeq\boldsymbol{0}$. Deuxièmement, la loss est supposée
localement bien décrite par son approximation quadratique, ce qui permet de
négliger les termes d'ordre supérieur. Troisièmement, les termes hors diagonale
de $H$ sont ignorés. Cette dernière approximation rend le calcul praticable,
mais elle néglige les compensations possibles entre paramètres, contrairement à
Optimal Brain Surgeon \cite{hassibi1993obs}.

Supprimer le poids $w_k$ revient à imposer
$\delta_k=-w_k$ et $\delta_j=0$ pour $j\neq k$. Sous les hypothèses
précédentes, l'augmentation locale de la loss devient
\begin{equation}
    \Delta\mathcal{L}_k
    \simeq \frac{1}{2}H_{kk}w_k^2.
\end{equation}
La saillance utilisée dans l'implémentation est donc
\begin{equation}
    S_k^{\mathrm{OBD}}
    =\frac{1}{2}\left|\widehat{H}_{kk}\right|w_k^2.
    \label{eq:score_obd}
\end{equation}
Dans la formulation idéale, la courbure diagonale est positive au voisinage
d'un minimum. La valeur absolue de $\widehat{H}_{kk}$ est une précaution de
l'implémentation face au bruit de l'estimation et à la non-convexité du réseau.
Un score faible signifie que l'annulation du poids devrait peu modifier la
loss, tandis qu'un score élevé signale un poids localement sensible.

\subsubsection{Estimation de la diagonale de la Hessienne}

Former explicitement la Hessienne complète demanderait une mémoire quadratique
en nombre de paramètres. Le code estime uniquement sa diagonale à partir de
produits Hessienne--vecteur. Pour un vecteur aléatoire de Rademacher
$\boldsymbol{v}\in\{-1,+1\}^{P}$ dont les coordonnées sont indépendantes,
\begin{equation}
    \mathbb{E}\!\left[
        \boldsymbol{v}\odot(H\boldsymbol{v})
    \right]
    =\operatorname{diag}(H).
\end{equation}
Avec $R$ tirages, l'estimateur employé est
\begin{equation}
    \widehat{\operatorname{diag}(H)}
    =\frac{1}{R}\sum_{r=1}^{R}
      \boldsymbol{v}^{(r)}\odot
      \left(H\boldsymbol{v}^{(r)}\right).
    \label{eq:hutchinson_diagonale}
\end{equation}
Chaque produit $H\boldsymbol{v}^{(r)}$ est obtenu par autodifférentiation sans
matérialiser $H$. Dans la configuration actuelle, la BCE est calculée sur 256
exemples d'entraînement et $R=10$ vecteurs sont utilisés. La diagonale est
estimée pour tous les paramètres entraînables, mais seuls les tenseurs dont le
nom se termine par \texttt{weight} participent au classement : les biais ne
sont pas prunés.

\subsubsection{Sélection globale et réentraînement itératif}

Les scores de toutes les matrices de poids sont aplatis et comparés dans un
classement global. Les poids déjà masqués reçoivent un score infini afin de ne
pas être sélectionnés une seconde fois. À chaque itération, le protocole retire
au plus $10\,\%$ du nombre initial de poids, puis réentraîne les poids actifs
pendant trois époques. Les gradients des poids masqués sont annulés et le masque
est réappliqué après chaque mise à jour, ce qui empêche les connexions supprimées
de réapparaître. Les scores et la diagonale de la Hessienne sont ensuite
recalculés sur le modèle réentraîné avant l'itération suivante.

La sparsité cible est définie par le rapport entre le nombre de poids de M1 et
celui de la référence dense M4 :
\begin{equation}
    \tau_H
    =1-\frac{P_{H,\mathrm{M4}}^{\mathrm{poids}}}
             {P_{H,\mathrm{M1}}^{\mathrm{poids}}}.
    \label{eq:sparsity_cible_m4}
\end{equation}
Elle vaut $50\,\%$ pour H1 et environ $49{,}91\,\%$ pour H2. La dernière
itération retire seulement le nombre de poids encore nécessaire pour atteindre
exactement cette cible.

\begin{table}[!ht]
    \centering
    \begin{tabularx}{\textwidth}{|l|X|}
        \hline
        Hyperparamètre OBD & Valeur actuelle \\ \hline
        Données du score & 256 premiers exemples de l'ensemble d'entraînement. \\ \hline
        Vecteurs de Rademacher & $R=10$ par estimation de la diagonale. \\ \hline
        Granularité & Poids individuels, classement global entre les couches. \\ \hline
        Pas d'élagage & $10\,\%$ du nombre initial de poids par itération, avec un dernier pas tronqué. \\ \hline
        Sparsité cible & H1 : $50\,\%$ ; H2 : $49{,}973\,\%$. \\ \hline
        Réentraînement intermédiaire & 3 époques, batch 256, Adam, taux $10^{-3}$, sans arrêt anticipé. \\ \hline
        Paramètres exclus & Biais et backbone Word2Vec. \\ \hline
    \end{tabularx}
    \caption{Configuration actuelle du pruning OBD.}
    \label{tab:hyperparametres_obd}
\end{table}

\subsection{Construction du score MINT}
\label{ssec:methode_mint}

\subsubsection{Dépendance conditionnelle entre deux couches}

MINT ne calcule son score ni à partir de la magnitude des poids, ni à partir de
leur gradient ou de leur Hessienne. La méthode étudie les activations produites
par deux couches successives sur un ensemble de calibration
\cite{ganesh2020mint}. Pour $N_{\mathrm{cal}}$ observations, les activations de
la couche $\ell$ sont rassemblées dans
\begin{equation}
    A^{(\ell)}\in\mathbb{R}^{N_{\mathrm{cal}}\times d_{\ell}},
\end{equation}
où une ligne correspond à un texte et une colonne à une unité du réseau. Les
sorties cachées sont observées après ReLU ; pour la dernière couche, le logit
est utilisé avant la sigmoïde.

Les $d_{\ell}$ unités amont sont réparties en $K_{\ell}$ groupes disjoints
$G_{\ell,1},\ldots,G_{\ell,K_{\ell}}$. Pour mesurer l'importance du groupe
$G_{\ell,g}$ vis-à-vis de l'unité aval $i$, trois variables sont construites :
\begin{align}
    X &= A^{(\ell+1)}_{:,i},
        & &\text{activation de l'unité aval},\\
    Y &= A^{(\ell)}_{:,G_{\ell,g}},
        & &\text{activations du groupe amont testé},\\
    Z &= A^{(\ell)}_{:,-G_{\ell,g}},
        & &\text{activations de toutes les autres unités amont}.
    \label{eq:triplet_mint}
\end{align}
Le score recherché est l'information mutuelle géométrique conditionnelle
\begin{equation}
    \rho_{i,g}^{(\ell)}
    =I_{\mathrm{G}}(X;Y\mid Z).
    \label{eq:score_mint_abstrait}
\end{equation}
Le conditionnement est essentiel. Un groupe $Y$ peut être corrélé à la sortie
$X$ tout en étant redondant si la même information est déjà contenue dans
$Z$. Un faible $\rho_{i,g}^{(\ell)}$ signifie donc que le groupe apporte peu
d'information supplémentaire à l'unité aval une fois les autres activations
connues.

Pour une observation, la concaténation $(X,Y,Z)$ appartient toujours à
$\mathbb{R}^{d_{\ell}+1}$, car
$1+|G_{\ell,g}|+(d_{\ell}-|G_{\ell,g}|)=d_{\ell}+1$. Par exemple, sur la
transition $1\,024\rightarrow128$, un groupe de 64 unités donne un scalaire
$X$, un vecteur $Y\in\mathbb{R}^{64}$ et un vecteur
$Z\in\mathbb{R}^{960}$ ; chaque point géométrique appartient donc à
$\mathbb{R}^{1\,025}$. Il faut plusieurs observations de ce triplet pour estimer
la dépendance : les coordonnées sont des activations, tandis que les lignes
représentent les textes de calibration.

\subsubsection{Du test à deux échantillons à l'indépendance conditionnelle}

La construction de l'estimateur se comprend d'abord comme un problème de
comparaison de deux distributions. Le \emph{Classifier Two-Sample Test}
(C2ST) fournit une intuition simple de ce principe
\cite{lopezpaz2017c2st}. Soient deux échantillons équilibrés
$S_P=\{U_r^P\}_{r=1}^{m}\sim P$ et
$S_Q=\{U_r^Q\}_{r=1}^{m}\sim Q$. Les observations de $S_P$ reçoivent le
label 1 et celles de $S_Q$ le label 0, puis l'ensemble
\begin{equation}
    \mathcal{D}
    =\{(U_r^P,1)\}_{r=1}^{m}
     \cup\{(U_r^Q,0)\}_{r=1}^{m}
\end{equation}
est séparé en un jeu d'entraînement et un jeu de test. Un classifieur binaire
$h:\mathcal{U}\rightarrow[0,1]$ est appris sur le premier jeu et son accuracy
sur des observations tenues à l'écart est mesurée par
\begin{equation}
    \widehat{t}
    =\frac{1}{|\mathcal{D}_{\mathrm{te}}|}
      \sum_{(u,c)\in\mathcal{D}_{\mathrm{te}}}
      \mathbf{1}\!\left[
        \mathbf{1}\!\left[h(u)>\frac{1}{2}\right]=c
      \right].
    \label{eq:c2st_accuracy}
\end{equation}
Si $P=Q$, le label de population n'est pas prédictible à partir de $u$ et
$\widehat{t}$ doit rester proche de $1/2$. Si $P\neq Q$ et si le classifieur
dispose d'une capacité suffisante, une accuracy supérieure au hasard révèle
que les deux populations sont séparables. Réciproquement, une accuracy proche
de $1/2$ ne prouve pas à elle seule l'égalité : elle conduit seulement à ne pas
rejeter cette hypothèse pour le classifieur et l'échantillon considérés. La
séparation entraînement--test est indispensable pour que la mémorisation ne
soit pas confondue avec une différence de distributions.

Pour tester l'indépendance ordinaire entre $X$ et $Y$, on compare la loi jointe
$P=P_{XY}$ à la loi produit $Q=P_XP_Y$. À partir des paires observées
$\{(X_r,Y_r)\}_{r=1}^{m}$, une permutation aléatoire $\pi$ des coordonnées
$Y$ fournit l'échantillon artificiel
$\{(X_r,Y_{\pi(r)})\}_{r=1}^{m}$ : l'association entre $X$ et $Y$ est cassée
tout en conservant leurs distributions marginales. Cette permutation globale
ne convient plus pour l'indépendance conditionnelle, car elle détruirait aussi
la relation entre $Y$ et $Z$. Il faut donc construire une population où $X$ et
$Y$ sont dissociés \emph{à $Z$ fixé}, au moins approximativement.

La GMI n'est pas l'information mutuelle de Shannon. Elle est reliée à une
divergence de Henze--Penrose et permet une estimation géométrique sans estimer
explicitement des densités en grande dimension \cite{yasaei2019gmi}. La
construction compare deux distributions définies sur les mêmes triplets :
\begin{align}
    P(x,y,z)
        &=f_{XYZ}(x,y,z),\\
    Q(x,y,z)
        &=f_Z(z)f_{X\mid Z}(x\mid z)f_{Y\mid Z}(y\mid z).
    \label{eq:distributions_mint_pq}
\end{align}
$P$ est la distribution réelle des activations. $Q$ est la distribution qui
serait obtenue si $X$ et $Y$ étaient rendus indépendants conditionnellement à
$Z$. Ainsi, $X\perp\!\!\!\perp Y\mid Z$ implique $P=Q$.

MINT conserve cette réduction à une comparaison entre $P$ et $Q$, mais
n'entraîne pas le classifieur du C2ST. L'accuracy de
l'\autoref{eq:c2st_accuracy} est remplacée par une statistique géométrique
calculée sur un arbre couvrant minimal. Le C2ST explique donc pourquoi rendre
les deux populations discernables mesure une dépendance ; le MST fournit le
score effectivement employé pour classer les blocs de poids.

\subsubsection{Population artificielle par bootstrap conditionnel}

Avant les calculs géométriques, chaque colonne d'activation est centrée et
réduite. Cette standardisation empêche les unités de grande variance de dominer
les distances euclidiennes. Pour chaque score, les observations sont ensuite
réparties en trois ensembles disjoints de taille $m$ :
\begin{equation}
    I_P,\ I_S,\ I_D,
    \qquad |I_P|=|I_S|=|I_D|=m.
\end{equation}
La population réelle contient les triplets
\begin{equation}
    \mathcal{P}
    =\{(X_r,Y_r,Z_r):r\in I_P\}.
\end{equation}

Une permutation globale de $Y$ casserait simultanément les relations
$X$--$Y$ et $Y$--$Z$ ; elle ne produirait donc pas la distribution $Q$. MINT
s'appuie plutôt sur le bootstrap au plus proche voisin de Sen et ses coauteurs
\cite{sen2017ci}. Pour chaque source $r\in I_S$, un donneur est recherché dans
$I_D$ uniquement à partir de la proximité de $Z$ :
\begin{equation}
    d^*(r)
    =\underset{s\in I_D}{\arg\min}\,
      \lVert Z_r-Z_s\rVert_2.
    \label{eq:donneur_mint}
\end{equation}
Le point synthétique conserve $(X_r,Z_r)$ mais emprunte $Y$ au donneur :
\begin{equation}
    \widetilde{U}_r=(X_r,Y_{d^*(r)},Z_r).
\end{equation}
La population
$\mathcal{Q}=\{\widetilde{U}_r:r\in I_S\}$ approxime ainsi la distribution
conditionnellement indépendante de l'\autoref{eq:distributions_mint_pq}.
Intuitivement, $(X_r,Z_r)$ apporte un tirage suivant
$f_Z(z)f_{X\mid Z}(x\mid z)$, tandis que le donneur proche de $Z_r$ fournit un
$Y$ approximativement distribué selon $f_{Y\mid Z}(y\mid z)$. Les ensembles
source et donneur sont disjoints pour ne pas recopier la paire réelle
$(X_r,Y_r)$.

À taille finie, la population générée n'est qu'une approximation. Sa qualité
dépend du volume de données et de la possibilité de trouver un donneur proche
dans l'espace de conditionnement. Lorsque $Z$ est très dimensionnel, les
distances entre voisins deviennent moins discriminantes. Cette limite est
particulièrement importante ici, puisque $Z$ peut contenir plusieurs centaines
d'activations.

\subsubsection{Estimateur géométrique par arbre couvrant minimal}

Les $m$ points réels et les $m$ points synthétiques sont fusionnés. Un arbre
couvrant minimal euclidien est construit sur ces $2m$ points, puis la
statistique de Friedman--Rafsky compte le nombre $R_m$ d'arêtes reliant une
observation réelle à une observation synthétique
\cite{friedman1979multivariate}. Plus précisément, si $T$ désigne le MST et
$c_r\in\{0,1\}$ le label de population attaché au sommet $r$, alors
\begin{equation}
    R_m
    =\sum_{\{r,s\}\in E(T)}
      \mathbf{1}\!\left[c_r\neq c_s\right].
    \label{eq:friedman_rafsky_mint}
\end{equation}
Les labels ne sont pas des coordonnées du MST : ils servent uniquement à
compter, après sa construction, les arêtes qui changent de population. Pour
deux populations de même taille,
l'estimateur utilisé par MINT se simplifie en
\begin{equation}
    \widehat{I}_{\mathrm{G}}(X;Y\mid Z)
    =1-\frac{R_m}{m}.
\end{equation}
L'implémentation borne cette quantité dans $[0,1]$ :
\begin{equation}
    \rho_{i,g}^{(\ell)}
    =\operatorname{clip}_{[0,1]}
      \left(1-\frac{R_m}{m}\right).
    \label{eq:score_mint_estime}
\end{equation}

L'interprétation est la suivante. Si $P$ et $Q$ sont proches, leurs points se
mélangent et le MST comporte de nombreuses arêtes croisées ; $R_m$ est grand et
le score est faible. Le remplacement de $Y$ modifie alors peu la distribution,
ce qui indique une contribution conditionnelle limitée. Si les deux
populations sont au contraire séparées, $R_m$ est faible et le score est élevé :
le groupe $Y$ contribue de manière importante à l'activation aval.
Le comportement attendu sous l'indépendance se résume donc par
\begin{equation}
    X\perp\!\!\!\perp Y\mid Z
    \quad\Longrightarrow\quad
    P=Q
    \quad\Longrightarrow\quad
    R_m\simeq m
    \quad\Longrightarrow\quad
    \widehat{I}_{\mathrm{G}}\simeq0.
\end{equation}

\subsubsection{Passage des scores aux blocs de poids}

Pour une transition $d_{\ell}\rightarrow d_{\ell+1}$, le calcul produit une
matrice
\begin{equation}
    \rho^{(\ell)}
    \in[0,1]^{d_{\ell+1}\times K_{\ell}}.
\end{equation}
Le score $(i,g)$ est associé au bloc de poids
$W^{(\ell+1)}_{i,G_{\ell,g}}$. Si ce score est sélectionné, tous les poids qui
relient le groupe $G_{\ell,g}$ à l'unité aval $i$ sont annulés. Avec des groupes
de taille supérieure à un, il s'agit donc d'un pruning par blocs dans chaque
ligne de la matrice, et non de la suppression d'un neurone complet.

L'article MINT formule la décision à partir d'un seuil $\delta$. Pour contrôler
la comparaison avec OBD, l'implémentation actuelle impose plutôt une sparsité
de blocs $\tau$ dans chaque couche et sélectionne les scores les plus faibles.
Une protection $\gamma$ limite la fraction maximale de blocs supprimés dans une
même couche :
\begin{equation}
    B_{\ell}^{\mathrm{prune}}
    =\min\!\left(
        \left\lfloor d_{\ell+1}K_{\ell}\tau\right\rfloor,
        \left\lfloor d_{\ell+1}K_{\ell}\gamma\right\rfloor
      \right).
    \label{eq:blocs_prunes_mint}
\end{equation}
Cette règle est une adaptation expérimentale : une même proportion de blocs
ne correspond exactement à une même proportion de poids que si tous les groupes
ont la même taille.

\subsubsection{Configuration appliquée à H2}

MINT est appliqué en une passe à H2. Une calibration équilibrée est tirée dans
l'ensemble d'entraînement avec 96 exemples de chaque classe, soit
$N_{\mathrm{cal}}=192$. Les labels servent uniquement à équilibrer ce tirage ;
ils n'entrent ni dans les triplets $(X,Y,Z)$, ni dans le MST. Chaque score
utilise trois partitions de $m=64$ observations, ce qui exploite exactement
$3m=192$ exemples.

\begin{table}[!ht]
    \centering
    \begin{tabularx}{\textwidth}{|l|r|r|r|r|X|}
        \hline
        Transition & $d_{\ell}$ & $d_{\ell+1}$ & $K_{\ell}$ & Taille & Scores \\ \hline
        $100\rightarrow1\,024$ & 100 & $1\,024$ & 10 & 10 & $1\,024\times10=10\,240$ \\ \hline
        $1\,024\rightarrow128$ & $1\,024$ & 128 & 16 & 64 & $128\times16=2\,048$ \\ \hline
        $128\rightarrow1$ & 128 & 1 & 8 & 16 & $1\times8=8$ \\ \hline
        Total & -- & -- & -- & -- & $12\,296$ estimations GMI \\ \hline
    \end{tabularx}
    \caption{Groupement des activations pour le scoring MINT de H2.}
    \label{tab:groupes_mint_h2}
\end{table}

La sparsité nominale de blocs est
$\tau=49{,}973\,\%$, valeur héritée de la comparaison H2/M1--H2/M4, et la
protection vaut $\gamma=80\,\%$. À cause des parties entières dans
l'\autoref{eq:blocs_prunes_mint}, les trois couches suppriment respectivement
$5\,117$, $1\,023$ et 3 blocs. Comme leurs tailles valent 10, 64 et 16, cela
représente $116\,690$ poids annulés sur $233\,600$, soit une sparsité
effective de $49{,}953\,\%$.

Le modèle immédiatement masqué constitue MINT-M2. Une copie est ensuite
réentraînée au plus dix époques avec Adam, un batch de 256, un taux
d'apprentissage de $10^{-3}$ et une patience de 4 ; le masque est maintenu après
chaque mise à jour. Le modèle obtenu constitue MINT-M3.

\begin{table}[!ht]
    \centering
    \begin{tabularx}{\textwidth}{|l|X|}
        \hline
        Hyperparamètre MINT & Valeur actuelle \\ \hline
        Modèle scoré & W1/H2/M1, architecture $100\rightarrow1\,024\rightarrow128\rightarrow1$. \\ \hline
        Calibration & 96 exemples par classe, tirés dans l'entraînement ; 192 au total. \\ \hline
        Standardisation & Colonne par colonne, avec $\varepsilon=10^{-8}$. \\ \hline
        Taille des partitions & $m=64$ pour la population réelle, les sources et les donneurs. \\ \hline
        Groupes par couche & $(K_0,K_1,K_2)=(10,16,8)$. \\ \hline
        Règle de sélection & Sparsité de blocs par couche $\tau=49{,}973\,\%$. \\ \hline
        Protection de couche & $\gamma=80\,\%$ au maximum. \\ \hline
        Graine & 42 pour la calibration ; 42, 43 et 44 pour les trois couches. \\ \hline
        Réentraînement final & Au plus 10 époques, batch 256, Adam, taux $10^{-3}$, patience 4. \\ \hline
        Paramètres exclus & Biais et backbone Word2Vec. \\ \hline
    \end{tabularx}
    \caption{Configuration actuelle du pruning MINT.}
    \label{tab:hyperparametres_mint}
\end{table}

\subsection{Témoin de pruning aléatoire}
\label{ssec:methode_random}

Nous ajoutons un pruning aléatoire comme point de comparaison. Il permet de
vérifier si OBD choisit vraiment des poids plus utiles qu'un simple tirage au
hasard. Pour une tête contenant $P_H$ poids éligibles et une sparsité cible
$\tau_H$, le nombre de poids supprimés est
\begin{equation}
    K_H=\operatorname{round}(P_H\tau_H).
    \label{eq:nombre_poids_random}
\end{equation}
Les matrices de poids sont d'abord aplaties dans un ordre fixe. Nous tirons
ensuite une permutation aléatoire uniforme $\pi$ de
$\{1,\ldots,P_H\}$. Le masque annule les indices
$\pi(1),\ldots,\pi(K_H)$. Tous les groupes de $K_H$ poids ont donc la même
chance d'être choisis. Ce tirage n'utilise ni les poids, ni les gradients, ni
les activations, ni les labels. Comme pour OBD, les biais et le backbone
Word2Vec ne sont pas prunés.

Pour comparer correctement cette méthode à OBD, nous gardons le même
calendrier. À chaque étape, le script supprime au plus $10\,\%$ des poids
initiaux, puis réentraîne les poids actifs pendant trois époques. H1 demande six
étapes, car la dernière ne supprime que deux poids à cause des arrondis. H2 en
demande cinq. RAND-M2 est le modèle obtenu juste après la dernière suppression,
avant le dernier réentraînement. Comme OBD-M2, il a déjà été réentraîné lors des
étapes précédentes. RAND-M3 est le modèle obtenu après le dernier
réentraînement.

L'expérience utilise les graines 42 à 46. Pour chaque graine, nous reconstruisons
M1 avec la même initialisation et le même ordre des mini-batches que dans les
expériences précédentes. Nous retrouvons exactement les mêmes métriques. Le
masque utilise une autre graine, obtenue en ajoutant $60\,000$ à la graine
principale. La sparsité finale est de $50\,\%$ pour H1 et de
$49{,}973\,\%$ pour H2. Le nombre de poids actifs est donc le même que pour
OBD.

\subsection{Protocole d'évaluation et précautions expérimentales}
\label{ssec:methode_evaluation}

La même fonction d'évaluation est utilisée pour toutes les méthodes. Elle
calcule la binary cross-entropy, l'accuracy, la balanced accuracy, la précision,
le rappel, le F1, le macro-F1 et la ROC-AUC sur les ensembles de validation et
de test. Le jeu de test n'intervient ni dans le calcul des scores, ni dans le
choix du masque, ni dans l'arrêt anticipé. Les métriques sont rapportées
pour M1, juste après le pruning avec M2, après réentraînement avec M3 et pour la
référence dense M4.

Les mesures de complexité distinguent le nombre de poids actifs, la taille
physique des checkpoints, les MAC théoriquement évitables et la latence
réellement mesurée. Cette distinction est nécessaire parce que les masques
actuels annulent des coefficients sans réduire les dimensions des tenseurs
denses. La mesure finale est effectuée sur CPU AMD Ryzen AI 9 365, sans CUDA,
avec 10 threads PyTorch. Elle utilise $2\,048$ textes de test préencodés, des
batches de 512, deux passages d'échauffement et la médiane de sept répétitions.
La tokenisation est exclue de la latence bout en bout rapportée.

Le run détaillé dans les figures repose sur la graine 42. Il est complété par
le protocole de répétitions décrit ci-après afin de vérifier si les écarts
observés survivent à la variabilité de l'optimisation et du calcul des scores.
Deux limites restent néanmoins présentes : le batch de scoring OBD est
constitué des 256 premiers exemples d'entraînement, et chaque score MINT repose
sur une seule permutation. Un tirage explicite du batch OBD et la mise en commun
de plusieurs partitions MINT permettraient une étude encore plus robuste.

\subsection{Répétitions appariées et estimation de la variabilité}
\label{ssec:methode_repetitions}

L'expérience complète est répétée avec les cinq graines 42, 43, 44, 45 et 46.
Le tokenizer, le backbone W1, les vecteurs agrégés et les partitions
entraînement--validation--test restent strictement identiques. Pour chaque
graine, les têtes M1 et M4 sont réinitialisées et réentraînées ; OBD répète
également les tirages de Rademacher utilisés pour estimer la diagonale de la
Hessienne, tandis que MINT répète le choix de l'ensemble de calibration et les
partitions employées par son estimateur géométrique. Une même graine identifie
donc un bloc apparié dans lequel toutes les méthodes voient les mêmes données
et le même backbone.

Pour chaque métrique $q$, les sorties conservent les cinq valeurs individuelles,
leur moyenne $\bar q$, leur variance empirique non biaisée
\begin{equation}
    s_q^2=\frac{1}{n-1}\sum_{r=1}^{n}(q_r-\bar q)^2,
    \qquad n=5,
\end{equation}
leur écart-type $s_q$ et un intervalle de confiance descriptif de la moyenne à
$95\,\%$, calculé avec la loi de Student à $n-1$ degrés de liberté. Les mêmes
statistiques sont calculées sur les différences graine par graine entre deux
modèles. Cette comparaison appariée est plus informative que deux intervalles
séparés, car elle retire une partie de la variabilité commune à la graine.

Ce protocole mesure la variabilité algorithmique conditionnellement au corpus,
au split et à W1. Il ne constitue ni une validation croisée, ni un rééchantillonnage
du jeu de test : les intervalles obtenus ne doivent donc pas être interprétés
comme l'incertitude totale de généralisation à une nouvelle population de
textes.
